\documentclass[a4paper, 11pt]{article}
\usepackage[utf8]{inputenc}
\usepackage[T1]{fontenc}
\usepackage[italian]{babel}
\usepackage{geometry}
\usepackage{graphicx}
\usepackage{tabularx}
\usepackage{fancyhdr}
\usepackage{hyperref}
\usepackage{xcolor}
\usepackage{enumitem}


% Margini
\geometry{top=3cm, bottom=3cm, left=2.5cm, right=2.5cm}

% Intestazione e piè di pagina
\pagestyle{fancy}
\fancyhf{}
\fancyhead[L]{\small \textit{0x06tudio} -- Preventivo Costi e Impegni Orari}
\fancyhead[R]{\small Anno Accademico 2026/2027}
\fancyfoot[C]{\thepage}
\renewcommand{\headrulewidth}{0.4pt}

\begin{document}
\thispagestyle{plain}

% --- INTESTAZIONE E LOGO ---
\begin{minipage}{0.6\textwidth}
    \vspace{-2cm}
    {\huge \textbf{0x06tudio}} \\[0.5em]
    {\large Progetto didattico di Ingegneria del Software} \\[0.3em]
    {\normalsize Università degli Studi di Padova}
\end{minipage}
\begin{minipage}{0.4\textwidth}
    \raggedleft
    \vspace{0pt}
    \includegraphics[width=6cm]{../../media/logo.png}
    \vspace{1cm}
    \small
    \href{mailto:0x06tudio@gmail.com}{0x06tudio@gmail.com}
\end{minipage}

\vspace{0.5cm}

\begin{center}
    {\LARGE \textbf{Preventivo Costi e Impegni Orari}} \\[0.3em]
    {\large \textbf{Capitolato C9: Zucchetti - Sviluppo di Skills per OpenCode per la verifica della tracciabilità end-
to-end da casi d'uso a requisiti e codice sorgente.}}
\end{center}

\vspace{0.5cm}

% --- INFORMAZIONI SUL DOCUMENTO ---
\begin{tabularx}{\textwidth}{l X}
    \textbf{Versione:} & 0.2.0 \\
    \textbf{Data:} & 2026/10/27 \\
    \textbf{Uso:} & Esterno / Contrattuale \\
    \textbf{Redattori:} & Alberto Cavalletto, Davide Colabove \\
    \textbf{Destinatari:} & Proff. Vardanega, Cardin, Zucchetti S.p.A. \\
\end{tabularx}

\vspace{1cm}

% --- REGISTRO DELLE MODIFICHE ---
\section*{Registro delle modifiche}
\begin{tabularx}{\textwidth}{|c|c|X|c|}
    \hline
    \textbf{Versione} & \textbf{Data} & \textbf{Descrizione} & \textbf{Autore / Ruolo} \\ \hline
    0.2.0 & 2026/10/27 & Correzioni finali e impaginazione in template & Davide Colabove \\ \hline
    0.1.0 & 2026/10/26 & Stesura iniziale della bozza del documento, ripartizione oraria dei ruoli e computo economico. & Alberto Cavalletto\\ \hline
\end{tabularx}

\newpage
\tableofcontents
\newpage

% --- 1. SCOPO DEL DOCUMENTO ---
\section{Scopo del documento}
Il presente documento ha lo scopo di illustrare e formalizzare la pianificazione delle risorse, degli impegni orari e dei relativi costi preventivati dal gruppo per lo sviluppo del progetto didattico relativo al \textbf{Capitolato C9}.

Nello specifico, vengono esposti:
\begin{itemize}[label=-]
    \item Il monte ore complessivo e la ripartizione del carico lavorativo tra i ruoli previsti dal corso;
    \item L'analisi delle mansioni assegnate a ciascun ruolo con specifica motivazione operativa calibrata sul Capitolato C9;
    \item La strategia concordata per la rotazione periodica dei membri nei ruoli;
    \item Il quadro economico riassuntivo calcolato in base alle tariffe orarie convenzionate;
    \item L'impegno ufficiale sulla data di consegna finale del prodotto.
\end{itemize}


% --- 2. IMPEGNI ORARI ---
\section{Impegni orari}
A seguito di un'approfondita analisi dei requisiti preliminari e della discussione interna tra i membri del gruppo, è stato stimato il fabbisogno orario complessivo necessario a garantire la massima qualità dei processi e dei prodotti intermedi.

Ciascuno dei 6 componenti del team si impegna a dedicare al progetto una quota individuale media di \textbf{92 ore produttive}, per un monte ore dell'organizzazione complessivo pari a \textbf{552 ore}.

Di seguito è riportata la suddivisione del monte ore preventivato per ciascun ruolo normato:

\vspace{0.4cm}

\begin{table}[h!]
\centering
\begin{tabularx}{\textwidth}{|l|c|c|X|}
    \hline
    \textbf{Ruolo} & \textbf{Ore Previste} & \textbf{Incidenza (\%)} & \textbf{Media per Membro (su 6)} \\ \hline
    Responsabile   & 48 h  & 8,70\%  & 8,0 h / componente \\ \hline
    Amministratore & 48 h  & 8,70\%  & 8,0 h / componente \\ \hline
    Analista       & 105 h & 19,02\% & 17,5 h / componente \\ \hline
    Progettista    & 125 h & 22,64\% & 20,8 h / componente \\ \hline
    Programmatore  & 136 h & 24,64\% & 22,7 h / componente \\ \hline
    Verificatore   & 90 h  & 16,30\% & 15,0 h / componente \\ \hline
    \textbf{TOTALE} & \textbf{552 h} & \textbf{100,00\%} & \textbf{92,0 h / componente} \\ \hline
\end{tabularx}
\caption{Ripartizione delle ore di lavoro produttivo tra i ruoli di progetto.}
\end{table}

\newpage

% --- 3. ANALISI DEI RUOLI ---
\section{Analisi dei ruoli e motivazioni operative per C9}

Le mansioni e l'allocazione oraria sono state ponderate in relazione alle sfide architetturali e metodologiche introdotte dal capitolato in questione:

\subsection{Responsabile (48 ore -- 8,70\%)}
Il Responsabile coordina l'avanzamento dei lavori, gestisce il calendario delle scadenze e pianifica le riunioni. Cura i rapporti formali e le comunicazioni con il Committente e il Proponente, approvando i verbali interni ed esterni. Il monte ore è calibrato per una supervisione costante lungo tutte le fasi del ciclo di vita.

\subsection{Amministratore (48 ore -- 8,70\%)}
L'Amministratore gestisce l'infrastruttura tecnologica e favorisce la collaborazione del gruppo. È responsabile della configurazione del repository Git, della predisposizione e manutenzione delle pipeline di CI/CD, della definizione degli ambienti di containerizzazione (Docker) e del supporto agli standard qualitativi della documentazione.

\subsection{Analista (105 ore -- 19,02\%)}
L'Analista formalizza l'Analisi dei Requisiti traducendo i bisogni del capitolato in casi d'uso, requisiti funzionali, prestazionali e di vincolo. Nel contesto del Capitolato C9, assume rilevanza la modellazione formale della tracciabilità mediante diagrammi (Mermaid, XMI e formati immagine quali .jpg e .PNG), richiedendo un monte ore significativo nelle fasi iniziali.

\subsection{Progettista (125 ore -- 22,64\%)}
Il Progettista definisce la struttura architetturale complessiva del software. Si occupa in modo specifico della definizione dell'architettura della skill OpenCode, della stesura dei contratti d'interfaccia dei tool deterministici e dell'ingegnerizzazione della dashboard grafica di sintesi.

\subsection{Programmatore (136 ore -- 24,64\%)}
Il Programmatore realizza l'implementazione del software secondo i canoni stabiliti dal Progettista. Il ruolo assorbe la quota oraria maggiore per consentire lo sviluppo modulare della skill OpenCode, l'implementazione dei tool deterministici in ambiente Node.js/JavaScript e l'interfacciamento reattivo della dashboard.

\subsection{Verificatore (90 ore -- 16,30\%)}
Il Verificatore presidia costantemente la qualità dei processi e degli artefatti. Effettua la verifica formale e sintattica della documentazione (AdR, NdP, ST) e si occupa della stesura, automazione ed esecuzione delle suite di test (unitari, di integrazione e di sistema).


% --- 4. ROTAZIONE DEI RUOLI ---
\section{Rotazione dei ruoli}
In linea con le buone pratiche dell'ingegneria del software e le linee guida del corso, il gruppo adotta un modello di \textbf{rotazione periodica a cadenza bisettimanale}, sincronizzato con gli sprint di sviluppo concordati.
\newline

Tale approccio risponde ai seguenti principi:
\begin{enumerate}
    \item \textbf{Crescita e consapevolezza trasversale:} Ogni componente del gruppo ricoprirà a turno ciascuno dei sei ruoli, sperimentando sia le responsabilità organizzative e di coordinamento, sia le sfide tecniche d'analisi, progettazione, codifica e verifica.
    \item \textbf{Indipendenza della verifica:} Viene applicata la rigida regola per cui nessun membro può assumere il ruolo di Verificatore per codice, diagrammi o sezioni documentali da lui stesso prodotti all'interno dello stesso ciclo di lavoro.
    \item \textbf{Continuità e transizione informativa:} Il passaggio di consegne avviene tramite retrospettive di sprint e aggiornamento dei repository, evitando colli di bottiglia o dipendenze critiche da singoli individui.
\end{enumerate}

\newpage

% --- 5. PREVENTIVO COSTI ---
\section{Preventivo costi}
La valorizzazione economica del progetto è calcolata applicando al computo delle ore le tariffe orarie definite dal tariffario didattico ufficiale per ciascun profilo professionale.

\vspace{0.4cm}

\begin{table}[h!]
\centering
\begin{tabularx}{\textwidth}{|l|c|c|c|X|}
    \hline
    \textbf{Ruolo} & \textbf{Ore Previste} & \textbf{Costo Orario} & \textbf{Totale Ruolo} & \textbf{Motivazione Operativa per C9} \\ \hline
    Responsabile   & 48 h  & € 30,00/h & € 1.440,00 & Gestione scadenze, verbali e coordinamento del gruppo. \\ \hline
    Amministratore & 48 h  & € 20,00/h & € 960,00   & Repository Git, pipeline CI/CD e ambienti di lavoro. \\ \hline
    Analista       & 105 h & € 25,00/h & € 2.625,00 & Formalizzazione requisiti e tracciabilità da diagrammi (Mermaid/XMI). \\ \hline
    Progettista    & 125 h & € 25,00/h & € 3.125,00 & Architettura skill OpenCode, contratti tool e struttura dashboard. \\ \hline
    Programmatore  & 136 h & € 15,00/h & € 2.040,00 & Sviluppo skill, tool deterministici in Node.js/JS e dashboard di sintesi. \\ \hline
    Verificatore   & 90 h  & € 15,00/h & € 1.350,00 & Verifica formale documenti (AdR, NdP, ST) e suite di test. \\ \hline
    \textbf{TOTALE} & \textbf{552 h} & \textbf{--} & \textbf{€ 11.540,00} & \textbf{Media: 92 ore/persona (6 componenti)} \\ \hline
\end{tabularx}
\caption{Quadro finanziario ufficiale del progetto per il Capitolato C9.}
\end{table}

\vspace{0.6cm}

\noindent
Il costo complessivo stimato per la realizzazione dell'intero capitolato ammonta pertanto a:
\begin{center}
    {\Large \textbf{€ 11.540,00} \quad (\textit{undicimilacinquecentoquaranta/00 euro})}
\end{center}

\vspace{1cm}

% --- 6. DATA DI CONSEGNA ---
\section{Data di consegna}
In considerazione dell'allocazione oraria concordata, dell'analisi di fattibilità e della sostenibilità economica del piano di lavoro, il gruppo si impegna a completare tutte le attività di sviluppo, verifica e collaudo relative al Capitolato C9 entro il giorno:

\vspace{0.6cm}

\begin{center}
    {\LARGE \textbf{30 Aprile 2027}}
\end{center}

\vspace{0.8cm}

\noindent
La pianificazione degli sprint e delle release intermedie consentirà il rilascio progressivo degli incrementi software e documentali, garantendo la piena rispondenza ai criteri di accettazione previsti dal committente.

\end{document}